\section{Financial Assessment and Investor Profiling}

The Assessment module collects user financial information and produces an individual financial profile that drives subsequent modules. The user supplies monthly income, monthly expenses, accessible savings, a primary financial goal, and optionally monthly EMI payments, whether EMI is already included in expenses, monthly savings capacity, a target amount, current goal savings and a requested goal timeline.

The assessment engine computes a bounded set of household indicators from these inputs using deterministic formulas. Monthly outflows are computed as monthly expenses plus EMI payments when EMI is not already included in expenses, ensuring EMI is never double-counted. Monthly surplus is income minus outflows. Savings rate is surplus divided by income, returning null when income is zero. Emergency fund coverage is accessible savings divided by monthly expenses, returning null when expenses are zero. Goal projection computes remaining amount, assumed monthly contribution, estimated months to goal, required contribution for a requested timeline and whether the goal is on track. All calculations retain full floating-point precision, with month rounding using a tolerance-aware ceiling helper that ignores binary floating-point noise.

The engine emits rule-based warnings and recommendations for negative surplus, zero surplus, zero income, zero expenses, zero accessible savings, coverage below three months, stated savings capacity exceeding surplus, emergency fund and debt repayment goals, and timeline overruns. The engine exposes its assumptions in the response, including that all amounts are monthly INR figures, that EMI handling depends on the inclusion flag, that no investment returns, taxes, inflation or interest are modelled, and that contributions are assumed constant.

Personalization is important in financial decision support because financial circumstances, obligations and objectives vary substantially between individuals. The assessment output forms the baseline for the recommendation engine and the Digital Twin. The recommendation engine derives risk capacity from income stability, positive surplus, emergency coverage and debt burden, and uses the lower of derived capacity and stated tolerance when matching category and fund risk bands. The Digital Twin uses the assessment as the baseline against which user-modified scenarios are compared.

The assessment is persisted for authenticated users through the assessments API with row-level security, and the latest assessment is surfaced on the home and finances pages for ongoing reference.